\documentclass[12pt]{article}
\usepackage[utf8]{inputenc}
\usepackage[T1]{fontenc}
\usepackage{hyperref}
\title{Tiang Language Resources for K-12 Curriculum Development}
\author{Compiled by Claude Sonnet 5 (claude-sonnet-5) for Andr\'as Kornai}
\date{}
\begin{document}
\maketitle
\section{Introduction}
Tiang (Djaul; ISO 639-3 tbj; Glottocode tian1237; Austronesian, Oceanic) is spoken on the eastern half of Dyaul Island, New Ireland Province, Papua New Guinea. Figures conflict: 790 (a 1972 report) versus roughly 4000 in newer sources. Documentation is sparse.
\section{Primary Educational Resources}
Christoph Holz (CQUniversity) produced bilingual Tiang--English story and picture books, described in a CQU news item. \url{https://www.cqu.edu.au/news/707661/passion-for-pini-drives-cqu-students-groundbreaking-djaul-island-publication}
\section{Linguistic and Reference Materials}
\begin{itemize}
\item \emph{A comprehensive grammar of Tiang} (Holz, CQUniversity dataset). \url{https://acquire.cqu.edu.au/articles/dataset/A_comprehensive_grammar_of_Tiang/23566965}
\item Researcher profile of Christoph Holz. \url{https://cqu.edu.au/research/current-research/profiles/598486/christoph-holz}
\item Tiang entry in the TLOPO database. \url{https://tlopo.clld.org/languages/tiang}
\item Wikipedia article. \url{https://en.wikipedia.org/wiki/Tiang_language}
\end{itemize}
\section{Religious and Community Resources}
None found.
\section{Digital and Audio Resources}
No audio found in search results (Wikipedia mentions archived materials, which I did not verify).
\section{Curriculum Development Context}
Speakers are multilingual and shifting toward Tok Pisin. The Holz grammar and story books are the most promising basis for school materials; collaboration with the author and the community would be needed.
\end{document}